\chapter*{Preface}
\addcontentsline{toc}{chapter}{Preface}

This book exists because of a gap that is easy to overlook.

The Islamic economic tradition contains a set of concepts ---
\arb{niyyah} (intention), \arb{qadar} (the determined), \arb{rizq}
(divinely apportioned provision), \arb{barakah} (the blessing that makes a
thing greater than its apparent magnitude), \arb{halal} and \arb{haram}
(the lawful and the forbidden), \arb{riba} (the prohibited interest) ---
that are not, in their classical formulation, claims about market
behavior. They are claims about the deep structure of reality and about
the relation between human action and that structure.

The contemporary literature on Islamic economics treats these concepts
in two ways. The first treatment reduces them to legal categories: a
\arb{halal} transaction is one that satisfies certain juristic
conditions; \arb{rizq} is what the agent receives by lawful means.
This treatment is necessary for practical jurisprudence but leaves the
metaphysical content unaddressed. The second treatment translates the
concepts into modern technical vocabularies: \arb{riba} becomes financial
entropy under non-ergodic dynamics; \arb{rizq} becomes a network
phenomenon under super-linear scaling; \arb{barakah} becomes a
multiplicative coefficient in a wealth equation. This treatment is
intellectually serious and has produced useful results, but it changes
what is being discussed. The classical concepts make claims about
\textit{being}; the modern translations make claims about
\textit{dynamics}. These are different layers of reality.

The work that has not been done, and that this book attempts to begin,
is the metaphysical treatment proper: an articulation of what these
classical concepts \textit{are} at the level at which they were
originally formulated, in a vocabulary that contemporary readers can
engage with rigorously. The vocabulary required is not the vocabulary of
applied mathematics. It is the vocabulary of metaphysics: the structure
of being, the relation between mind and world, the nature of
actualization from potentiality, the limits of computation, the
hiddenness of the deep order beneath the manifest one.

The work is built on a deliberate methodological discipline: the
apparatus drawn upon must have genuine scientific authority. It is not
enough that a framework be philosophically interesting or
phenomenologically suggestive; it must carry the kind of empirical or
mathematical weight that compels assent across different schools of
thought. This rules out a great deal of contemporary metaphysics ---
process philosophy, phenomenology, analytical idealism --- as
load-bearing apparatus. Each of those traditions has produced thinkers
whose work parallels Akbarian themes in suggestive ways. None of them,
however, has the standing to ground a metaphysical articulation of
classical Islamic doctrine in a way that engages a reader who comes from
a scientific worldview.

Two bodies of contemporary work do meet that standard, and together they
suffice. The first is the foundations of quantum mechanics: the
measurement problem, the live interpretations (Bohmian, QBist,
Many-Worlds), and the non-locality established beyond reasonable doubt
by the loophole-free Bell experiments. Quantum mechanics is not just
suggestive; it is the most successful theory in the history of science,
and it has forced the abandonment of the classical materialist picture.
That makes it the only major scientific framework that opens the
categorial space in which a two-layer ontology --- the relation between
the manifest and the Unseen --- can be articulated.

The second is mathematical logic and computability theory, specifically
the argument from G\"odel's incompleteness theorems and the Turing
halting problem to the non-algorithmic character of mathematical
understanding. This argument, developed most fully by Roger Penrose,
gives a precise mathematical sense in which understanding --- and
\textit{a fortiori} intention --- cannot be a computational state. The
authority here is not empirical but mathematical, and it is settled.
This second framework provides a lever for articulating \arb{niyyah}
that is independent of quantum mechanics and that does not depend on any
particular interpretation of physics.

These two scientific frameworks do not, by themselves, articulate
Islamic metaphysics. They open the categorial space. The articulation
is then done in the indigenous vocabulary of the Akbarian tradition ---
\arb{wujud}, \arb{tajalli}, the divine names, the structure of
self-disclosure as developed by Ibn \arb{Arabi}, Mulla \arb{Sadra}, and
Shah \arb{Wali Allah}. The book takes the contemporary scientific
apparatus and the classical Akbarian apparatus to be saying compatible
things about the structure of reality. Its aim is to make the
compatibility visible and to use it to articulate what the classical
Islamic economic concepts are claiming about reality.

A note on what this book deliberately omits. The applied
mathematical treatments of Islamic economic concepts that have been
developed over the past several decades --- ergodicity for \arb{riba},
network theory for family and community provision, behavioral mediation
for the effects of gratitude, debt-deflation theorems for collapse
dynamics --- are valuable. They produce testable predictions and engage
with mainstream economics in productive ways. They do not appear in
this book. They belong to a separate companion volume in heterodox
Islamic economics, where their methodological assumptions and their
relation to the metaphysical claims treated here can be addressed
properly. Mixing the two registers, as my own earlier draft did, blurs
both. Each project is better served by being itself.

A note on register. The book is written for readers who are willing to
take metaphysics seriously as a source of knowledge --- not as
decoration, and not as a placeholder for what science has not yet
explained, but as a discipline with its own subject matter and methods.
That subject matter is the structure of being. Its methods are
conceptual analysis, careful attention to texts, and the use of formal
apparatus where formal apparatus genuinely illuminates what is at issue.
The book assumes some economic literacy but builds the required physics,
philosophy, and Akbarian background from foundations.

A note on the relation to subsequent work. This book is the
metaphysical foundation. A subsequent volume (Volume I in the planned
sequence) will build the formal microeconomic theory --- preferences,
choice, the firm, equilibrium, welfare --- on this foundation, in the
mode of Mas-Colell, Whinston, and Green but with the Islamic ontology
made explicit at every step. A separate companion in heterodox economics
will take up the applied mathematical treatments. The metaphysical
work, however, is prior to both. Without it, the formal theory has no
ground that distinguishes it from any other formal theory; and without
it, the heterodox treatments are at risk of mistaking a useful model for
a description of what the modeled phenomenon is.

\vspace{1em}
\noindent\textit{Working draft. Comments and corrections sought.}
